\documentclass[10pt]{extarticle}
\usepackage[a4paper,margin=1.5cm]{geometry}
\usepackage[T1]{fontenc}
\usepackage{amsmath,amssymb}
\usepackage{graphicx}
\usepackage[export]{adjustbox}
\usepackage{multicol}
\usepackage{enumitem}
\usepackage{booktabs,array,tabularx}
\setlength{\columnsep}{0.7cm}
\setlength{\parindent}{0pt}
\setlength{\parskip}{3pt}
\pagestyle{empty}
\makeatletter
% figure registry (figdefs.tex is generated by the extract step)
\newcommand{\qfigdef}[4]{\@namedef{qf@file@#1}{#2}\@namedef{qf@w@#1}{#3}\@namedef{qf@wide@#1}{#4}}
\newcommand{\qf@check}[1]{\@ifundefined{qf@file@#1}{\PackageError{score2}{Unknown figure id #1}{}}{}}
\newcommand{\qf@img}[4][]{\includegraphics[width=#2\dimexpr\@nameuse{qf@w@#4}pt\relax,max width=#3,max height=0.45\textheight,#1]{\@nameuse{qf@file@#4}}}
% Full-width questions: finalize wraps every question that contains a wide figure in \qfullbegin...\qfullend,
% so the WHOLE question (stem, figures, options) is set across the page in reading order, then the two columns
% resume. The question's head - its number, stem and first figure (\qfullsplit marks the end of the head) -
% is kept on one page, so a number is never left at a page foot with its figure overleaf; the rest may break
% across pages normally, so no big gaps are left. Without a \qfullsplit, the whole question is kept together
% if it is under 45% of a page. Never write these yourself.
\newif\ifqf@full
\newif\ifqf@boxed
\newsavebox{\qf@box}
\newcommand{\qf@open}{\global\qf@fulltrue\par\global\qf@boxedtrue
  \setbox\qf@box\vbox\bgroup\hsize\textwidth\linewidth\textwidth}
\newcommand{\qf@emit}[1]{\par\egroup\global\qf@boxedfalse
  \ifdim\dimexpr\ht\qf@box+\dp\qf@box\relax<#1\textheight\noindent\usebox{\qf@box}\par\else\unvbox\qf@box\fi}
\newcommand{\qfullsplit}{\ifqf@boxed\qf@emit{0.85}\fi}
\newcommand{\qf@close}{\ifqf@boxed\qf@emit{0.45}\fi\global\qf@fullfalse}
\newcommand{\qfullbegin}{\end{multicols}\qf@open}
\newcommand{\qfullend}{\qf@close\begin{multicols}{2}}
\newcommand{\qfullstart}{\qf@open}
\newcommand{\qfulldone}{\qf@close}
% \qfig[scale]{id}: block figure, centred, directly where it is written (size set by the tool).
% Fallback if a wide figure is not inside a full-width question: break out of the columns in place - no floats.
\newcommand{\qfig}[2][1]{\qf@check{#2}%
  \ifqf@full\qf@blk{#1}{#2}\else
  \ifnum\@nameuse{qf@wide@#2}=1 \expandafter\@firstoftwo\else\expandafter\@secondoftwo\fi
  {\end{multicols}\begin{center}\qf@img{#1}{\textwidth}{#2}\end{center}\begin{multicols}{2}}%
  {\qf@blk{#1}{#2}}\fi}
\newcommand{\qf@blk}[2]{\par\smallskip{\centering\qf@img{#1}{\linewidth}{#2}\par}\smallskip}
% \qfigi[scale]{id}: inline figure (inside option lists / tables / minipages); never breaks columns.
\newcommand{\qfigi}[2][1]{\qf@check{#2}\qf@img[valign=c]{#1}{\linewidth}{#2}}
\makeatother
% structure macros - lint reads these to check completeness and figure placement
\newcommand{\papersec}[1]{\par\bigskip{\centering\Large\bfseries PAPER-#1\par}\medskip}
\newcommand{\partsec}[2]{\par\medskip\noindent{\large\bfseries PART-#1: #2}\par\smallskip}
\newcommand{\sect}[1]{\par\smallskip\noindent\textbf{#1}\par}
\newcommand{\qn}[1]{\par\smallskip\noindent\textbf{#1.}~}
\newcommand{\paraq}[2]{\par\smallskip\noindent\textbf{Paragraph for Questions #1 and #2}\par}
\newenvironment{qoptions}{\begin{itemize}[leftmargin=1.8em,itemsep=1pt,topsep=2pt,label={}]}{\end{itemize}}
\newcommand{\opt}[2]{\item[(#1)] #2}
% figure registry (generated by score2_tools.py finalize)
\qfigdef{p1_phy_s1_q01_1}{Score Advanced - Test 27 (24 Apr 2024)-img/p1_phy_s1_q01_1.png}{78.2}{0}
\qfigdef{p1_phy_s1_q02_1}{Score Advanced - Test 27 (24 Apr 2024)-img/p1_phy_s1_q02_1.png}{108.2}{0}
\qfigdef{p1_phy_s1_q03_1}{Score Advanced - Test 27 (24 Apr 2024)-img/p1_phy_s1_q03_1.png}{118.9}{0}
\qfigdef{p1_phy_s1_q04_1}{Score Advanced - Test 27 (24 Apr 2024)-img/p1_phy_s1_q04_1.png}{110.7}{0}
\qfigdef{p1_phy_s2_q05_1}{Score Advanced - Test 27 (24 Apr 2024)-img/p1_phy_s2_q05_1.png}{61.2}{0}
\qfigdef{p1_phy_s2_q07_1}{Score Advanced - Test 27 (24 Apr 2024)-img/p1_phy_s2_q07_1.png}{139.3}{0}
\qfigdef{p1_phy_s2_q08_1}{Score Advanced - Test 27 (24 Apr 2024)-img/p1_phy_s2_q08_1.png}{105.6}{0}
\qfigdef{p1_phy_s2_q09_1}{Score Advanced - Test 27 (24 Apr 2024)-img/p1_phy_s2_q09_1.png}{130.0}{0}
\qfigdef{p1_phy_s2_q11_1}{Score Advanced - Test 27 (24 Apr 2024)-img/p1_phy_s2_q11_1.png}{201.1}{0}
\qfigdef{p1_phy_s2_q12_1}{Score Advanced - Test 27 (24 Apr 2024)-img/p1_phy_s2_q12_1.png}{209.9}{0}
\qfigdef{p1_phy_s3_q02_1}{Score Advanced - Test 27 (24 Apr 2024)-img/p1_phy_s3_q02_1.png}{104.9}{0}
\qfigdef{p1_phy_s3_q06_1}{Score Advanced - Test 27 (24 Apr 2024)-img/p1_phy_s3_q06_1.png}{108.0}{0}
\qfigdef{p1_chem_s1_q01_A}{Score Advanced - Test 27 (24 Apr 2024)-img/p1_chem_s1_q01_A.png}{149.4}{0}
\qfigdef{p1_chem_s1_q01_B}{Score Advanced - Test 27 (24 Apr 2024)-img/p1_chem_s1_q01_B.png}{175.3}{0}
\qfigdef{p1_chem_s1_q01_C}{Score Advanced - Test 27 (24 Apr 2024)-img/p1_chem_s1_q01_C.png}{180.1}{0}
\qfigdef{p1_chem_s1_q01_D}{Score Advanced - Test 27 (24 Apr 2024)-img/p1_chem_s1_q01_D.png}{163.0}{0}
\qfigdef{p1_chem_s1_q02_1}{Score Advanced - Test 27 (24 Apr 2024)-img/p1_chem_s1_q02_1.png}{42.1}{0}
\qfigdef{p1_chem_s1_q02_A}{Score Advanced - Test 27 (24 Apr 2024)-img/p1_chem_s1_q02_A.png}{31.3}{0}
\qfigdef{p1_chem_s1_q02_B}{Score Advanced - Test 27 (24 Apr 2024)-img/p1_chem_s1_q02_B.png}{46.5}{0}
\qfigdef{p1_chem_s1_q02_C}{Score Advanced - Test 27 (24 Apr 2024)-img/p1_chem_s1_q02_C.png}{33.8}{0}
\qfigdef{p1_chem_s1_q02_D}{Score Advanced - Test 27 (24 Apr 2024)-img/p1_chem_s1_q02_D.png}{33.7}{0}
\qfigdef{p1_chem_s2_q05_A}{Score Advanced - Test 27 (24 Apr 2024)-img/p1_chem_s2_q05_A.png}{174.9}{0}
\qfigdef{p1_chem_s2_q07_1}{Score Advanced - Test 27 (24 Apr 2024)-img/p1_chem_s2_q07_1.png}{211.8}{0}
\qfigdef{p1_chem_s2_q07_A}{Score Advanced - Test 27 (24 Apr 2024)-img/p1_chem_s2_q07_A.png}{31.4}{0}
\qfigdef{p1_chem_s2_q07_B}{Score Advanced - Test 27 (24 Apr 2024)-img/p1_chem_s2_q07_B.png}{53.1}{0}
\qfigdef{p1_chem_s2_q07_C}{Score Advanced - Test 27 (24 Apr 2024)-img/p1_chem_s2_q07_C.png}{74.5}{0}
\qfigdef{p1_chem_s2_q07_D}{Score Advanced - Test 27 (24 Apr 2024)-img/p1_chem_s2_q07_D.png}{28.4}{0}
\qfigdef{p1_chem_s2_q08_1}{Score Advanced - Test 27 (24 Apr 2024)-img/p1_chem_s2_q08_1.png}{245.2}{0}
\qfigdef{p1_chem_s2_q12_A}{Score Advanced - Test 27 (24 Apr 2024)-img/p1_chem_s2_q12_A.png}{104.8}{0}
\qfigdef{p1_chem_s2_q12_B}{Score Advanced - Test 27 (24 Apr 2024)-img/p1_chem_s2_q12_B.png}{133.4}{0}
\qfigdef{p1_chem_s2_q12_C}{Score Advanced - Test 27 (24 Apr 2024)-img/p1_chem_s2_q12_C.png}{130.8}{0}
\qfigdef{p1_chem_s2_q12_D}{Score Advanced - Test 27 (24 Apr 2024)-img/p1_chem_s2_q12_D.png}{124.9}{0}
\qfigdef{p2_phy_s1_q01_1}{Score Advanced - Test 27 (24 Apr 2024)-img/p2_phy_s1_q01_1.png}{131.5}{0}
\qfigdef{p2_phy_s1_q02_1}{Score Advanced - Test 27 (24 Apr 2024)-img/p2_phy_s1_q02_1.png}{131.7}{0}
\qfigdef{p2_phy_s1_q04_1}{Score Advanced - Test 27 (24 Apr 2024)-img/p2_phy_s1_q04_1.png}{106.1}{0}
\qfigdef{p2_phy_s1_q05_1}{Score Advanced - Test 27 (24 Apr 2024)-img/p2_phy_s1_q05_1.png}{128.3}{0}
\qfigdef{p2_phy_s1_q06_1}{Score Advanced - Test 27 (24 Apr 2024)-img/p2_phy_s1_q06_1.png}{124.6}{0}
\qfigdef{p2_phy_s1_q08_1}{Score Advanced - Test 27 (24 Apr 2024)-img/p2_phy_s1_q08_1.png}{138.9}{0}
\qfigdef{p2_phy_s2_para9_10_1}{Score Advanced - Test 27 (24 Apr 2024)-img/p2_phy_s2_para9_10_1.png}{52.5}{0}
\qfigdef{p2_phy_s2_para9_10_2}{Score Advanced - Test 27 (24 Apr 2024)-img/p2_phy_s2_para9_10_2.png}{56.8}{0}
\qfigdef{p2_phy_s2_para9_10_3}{Score Advanced - Test 27 (24 Apr 2024)-img/p2_phy_s2_para9_10_3.png}{45.1}{0}
\qfigdef{p2_phy_s2_para9_10_4}{Score Advanced - Test 27 (24 Apr 2024)-img/p2_phy_s2_para9_10_4.png}{39.2}{0}
\qfigdef{p2_phy_s2_para9_10_5}{Score Advanced - Test 27 (24 Apr 2024)-img/p2_phy_s2_para9_10_5.png}{48.2}{0}
\qfigdef{p2_phy_s2_para11_12_1}{Score Advanced - Test 27 (24 Apr 2024)-img/p2_phy_s2_para11_12_1.png}{93.0}{0}
\qfigdef{p2_phy_s2_para11_12_III}{Score Advanced - Test 27 (24 Apr 2024)-img/p2_phy_s2_para11_12_III.png}{78.7}{0}
\qfigdef{p2_phy_s2_para11_12_IV}{Score Advanced - Test 27 (24 Apr 2024)-img/p2_phy_s2_para11_12_IV.png}{82.8}{0}
\qfigdef{p2_phy_s3_q04_1}{Score Advanced - Test 27 (24 Apr 2024)-img/p2_phy_s3_q04_1.png}{114.8}{0}
\qfigdef{p2_chem_s1_q01_D}{Score Advanced - Test 27 (24 Apr 2024)-img/p2_chem_s1_q01_D.png}{28.7}{0}
\qfigdef{p2_chem_s1_q02_A}{Score Advanced - Test 27 (24 Apr 2024)-img/p2_chem_s1_q02_A.png}{105.8}{0}
\qfigdef{p2_chem_s1_q02_B}{Score Advanced - Test 27 (24 Apr 2024)-img/p2_chem_s1_q02_B.png}{101.1}{0}
\qfigdef{p2_chem_s1_q02_C}{Score Advanced - Test 27 (24 Apr 2024)-img/p2_chem_s1_q02_C.png}{81.1}{0}
\qfigdef{p2_chem_s1_q02_D}{Score Advanced - Test 27 (24 Apr 2024)-img/p2_chem_s1_q02_D.png}{119.0}{0}
\qfigdef{p2_chem_s1_q03_A}{Score Advanced - Test 27 (24 Apr 2024)-img/p2_chem_s1_q03_A.png}{121.6}{0}
\qfigdef{p2_chem_s1_q03_B}{Score Advanced - Test 27 (24 Apr 2024)-img/p2_chem_s1_q03_B.png}{135.6}{0}
\qfigdef{p2_chem_s1_q03_C}{Score Advanced - Test 27 (24 Apr 2024)-img/p2_chem_s1_q03_C.png}{116.4}{0}
\qfigdef{p2_chem_s1_q03_D}{Score Advanced - Test 27 (24 Apr 2024)-img/p2_chem_s1_q03_D.png}{101.2}{0}
\qfigdef{p2_chem_s2_para9_10_1}{Score Advanced - Test 27 (24 Apr 2024)-img/p2_chem_s2_para9_10_1.png}{112.2}{0}
\qfigdef{p2_chem_s2_para9_10_5}{Score Advanced - Test 27 (24 Apr 2024)-img/p2_chem_s2_para9_10_5.png}{23.5}{0}
\qfigdef{p2_chem_s2_para9_10_2}{Score Advanced - Test 27 (24 Apr 2024)-img/p2_chem_s2_para9_10_2.png}{106.6}{0}
\qfigdef{p2_chem_s2_para9_10_6}{Score Advanced - Test 27 (24 Apr 2024)-img/p2_chem_s2_para9_10_6.png}{23.2}{0}
\qfigdef{p2_chem_s2_para9_10_3}{Score Advanced - Test 27 (24 Apr 2024)-img/p2_chem_s2_para9_10_3.png}{142.6}{0}
\qfigdef{p2_chem_s2_para9_10_7}{Score Advanced - Test 27 (24 Apr 2024)-img/p2_chem_s2_para9_10_7.png}{85.3}{0}
\qfigdef{p2_chem_s2_para9_10_4}{Score Advanced - Test 27 (24 Apr 2024)-img/p2_chem_s2_para9_10_4.png}{110.1}{0}
\qfigdef{p2_chem_s2_para9_10_8}{Score Advanced - Test 27 (24 Apr 2024)-img/p2_chem_s2_para9_10_8.png}{73.2}{0}
\qfigdef{p2_chem_s2_para9_10_9}{Score Advanced - Test 27 (24 Apr 2024)-img/p2_chem_s2_para9_10_9.png}{61.7}{0}
\qfigdef{p2_chem_s2_para9_10_10}{Score Advanced - Test 27 (24 Apr 2024)-img/p2_chem_s2_para9_10_10.png}{63.4}{0}
\qfigdef{p2_math_s1_q01_1}{Score Advanced - Test 27 (24 Apr 2024)-img/p2_math_s1_q01_1.png}{132.5}{0}
\raggedcolumns
\makeatletter\@beginparpenalty=10000\makeatother
\newcolumntype{L}[1]{>{\raggedright\arraybackslash}p{#1}}
\begin{document}
\begin{center}{\Large\bfseries Score Advanced --- Test 27 (24 Apr 2024)}\end{center}
\begin{multicols}{2}
\raggedcolumns
\papersec{1}
\partsec{1}{PHYSICS}
\sect{SECTION-I (i)}
\qn{1} Consider the earth as a uniform sphere of mass m and radius R. Imagine a straight smooth tunnel made through the earth which connects any two points on its surface as shown in the figure. The particle is released from rest as shown in the figure. Find the time when it comes at the same point (g = GM / R$^2$)\par
\qfig{p1_phy_s1_q01_1}
\begin{qoptions}
\opt{A}{$2\pi\sqrt{\dfrac{R_e}{g}}$}
\opt{B}{$2\pi\sqrt{\dfrac{3R_e}{4g}}$}
\opt{C}{$2\pi\sqrt{\dfrac{R_e}{2g}}$}
\opt{D}{$2\pi\sqrt{\dfrac{3R_e}{g}}$}
\end{qoptions}
\qn{2} Two right angled prisms having the same refracting angle 60$^\circ$ are placed as shown in the figure. A ray of light incident on the first prism is finally deviated by an angles 90$^\circ$ grazingly as shown. The refractive index of second prism is $\sqrt2$. The refractive index of the first prism is equal to\par
\qfig{p1_phy_s1_q02_1}
\begin{qoptions}
\opt{A}{$\sqrt2\,\dfrac{\sin30^\circ}{\sin60^\circ}$}
\opt{B}{$\sqrt2\,\dfrac{\sin45^\circ}{\sin60^\circ}$}
\opt{C}{$\sqrt2\,\dfrac{\sin15^\circ}{\sin60^\circ}$}
\opt{D}{$\sqrt2\,\dfrac{\sin60^\circ}{\sin30^\circ}$}
\end{qoptions}
\qn{3} A hemisphere of radius R and density $\rho$ is held in space as shown in figure. If gravitational field intensity at P is E$(-\hat{\mathrm j})$ then magnitude of gravitational field intensity at Q is\par
\qfig{p1_phy_s1_q03_1}
\begin{qoptions}
\opt{A}{$\left(\mathrm E-\dfrac{\pi \mathrm{GR}\rho}{3}\right)(\hat{\mathrm j})$}
\opt{B}{$\dfrac{\pi \mathrm{GR}\rho}{3}(-\hat{\mathrm j})$}
\opt{C}{$\left(\mathrm E-\dfrac{\pi \mathrm{GR}\rho}{3}\right)(-\hat{\mathrm j})$}
\opt{D}{$\left(\mathrm E+\dfrac{\pi \mathrm{GR}\rho}{3}\right)(\hat{\mathrm j})$}
\end{qoptions}
\qn{4} Find the value of r so that maximum power is generated in external circuit.\par
\qfig{p1_phy_s1_q04_1}
\begin{qoptions}
\opt{A}{r = 10$\Omega$}
\opt{B}{r = 20$\Omega$}
\opt{C}{r = 30$\Omega$}
\opt{D}{r = 40$\Omega$}
\end{qoptions}
\sect{SECTION-I (ii)}
\qn{5} A spherical shell of volume mass density $\rho$, thickness t and radius R (R$\gg$t) is placed concentrically inside another shell of radius 2R having same thickness and of same material as shown in figure. The circumferential stress developed in the outer shell is\par
\qfig{p1_phy_s2_q05_1}
\begin{qoptions}
\opt{A}{Compressive in nature}
\opt{B}{Tensile in nature}
\opt{C}{$3\mathrm{G}\pi \mathrm{Rt}\rho^2$}
\opt{D}{$\dfrac32\mathrm{G}\pi \mathrm{Rt}\rho^2$}
\end{qoptions}
\qn{6} An $_{8}\mathrm{O}^{16}$ nucleus is spherical having radius R, surface area A and volume V. According to empirical observations, the volume of the $_{50}\mathrm{Sn}^{128}$ assumed to be spherical is V$'$, radius R$'$ and surface area A$'$. Then
\begin{qoptions}
\opt{A}{R$'$ = 8R}
\opt{B}{A$'$ = 8A}
\opt{C}{V$'$ = 8V}
\opt{D}{R$'$ = 2R}
\end{qoptions}
\qn{7} Two objects of the same mass slide down on two different smooth fixed inclined plane of same height but of different angles. When they reach the bottom, then I$_1$ and I$_2$ are the impulses of the weight of the first and the second object during the motion, P$_1$ and P$_2$ are the powers developed by the weight of the first and second objects.\par
\qfig{p1_phy_s2_q07_1}
\begin{qoptions}
\opt{A}{$\dfrac{\mathrm I_1}{\mathrm I_2}=\dfrac{\sin\alpha_2}{\sin\alpha_1}$}
\opt{B}{$\dfrac{\mathrm I_1}{\mathrm I_2}=1$}
\opt{C}{$\dfrac{\mathrm I_1}{\mathrm I_2}=\dfrac{\sin\alpha_1}{\sin\alpha_2}$}
\opt{D}{$\dfrac{\mathrm P_1}{\mathrm P_2}=\dfrac{\sin\alpha_1}{\sin\alpha_2}$}
\end{qoptions}
\qn{8} The ends of a long bar are maintained at different temperature and there is no lose of heat from the sides of the bar due to conduction or radiation. The graph of temperature against distance from the hot end of the bar at steady state is shown in the diagram. From this we can conclude that\par
\qfig{p1_phy_s2_q08_1}
\begin{qoptions}
\opt{A}{The cross-sectional area of the bar may be increasing as that distance from the hot end increases}
\opt{B}{The cross-sectional area of the bar may be decreasing as the distance from the hot end increases}
\opt{C}{The thermal conductivity of the bar may be decreasing as the distance from the hot end increases}
\opt{D}{The thermal conductivity of the bar may be increasing as that distance from the hot end increase}
\end{qoptions}
\qn{9} Consider the circuit shown below. I$_\mathrm S$ is constant current source, meaning that no matter what device is connected between points A and B, the current provided by the constant current source is always (I$_\mathrm S$) the same\par
\qfig{p1_phy_s2_q09_1}
\begin{qoptions}
\opt{A}{If an ideal voltmeter is connected between A and B, the reading of voltmeter is I$_\mathrm S$R.}
\opt{B}{If an ideal voltmeter is connected between A and B, the reading of volume is 2I$_\mathrm S$R.}
\opt{C}{If an ideal ammeter is connected between A and B, the reading of ammeter is $\dfrac{\mathrm I_S}{9}$.}
\opt{D}{If an ideal ammeter is connected between A and B, the reading of ammeter is $\dfrac{\mathrm I_S}{6}$.}
\end{qoptions}
\qn{10} We have two equal size boxes A and B. Each box contains gas that behaves as an ideal gas. Thermometer are inserted into each box and it is found that the gas in box A is at a temperature of 50$^\circ$C, while the gas in box B is at 10$^\circ$C. Which of the following statements(s) is/are true?
\begin{qoptions}
\opt{A}{The pressure in A may be higher than in B.}
\opt{B}{There may be more molecules in A and than in B.}
\opt{C}{The molecules in A have more average kinetic energy per molecule than those in B.}
\opt{D}{The molecules in A are moving faster than B.}
\end{qoptions}
\qn{11} A helicopter has a mass m and maintains height by imparting a downward momentum to a column of air defined by the slipstream boundary as shown in figure. The propeller blades can project a downward air speed v, where the pressure in the stream below the blade is atmospheric and the radius of the circular cross section of the slip stream is r. Neglect any rotational energy of the air, the temperature rise due to air friction and any change in air density ($\rho$)\par
\qfig{p1_phy_s2_q11_1}
\begin{qoptions}
\opt{A}{V $=\dfrac1r\sqrt{\dfrac{\mathrm{mg}}{\pi\rho}}$ for equilibrium}
\opt{B}{$\nu=\dfrac2r\sqrt{\dfrac{\mathrm{mg}}{\pi\rho}}$ for equilibrium}
\opt{C}{Power of engine $=\dfrac{\mathrm{mg}}{2r}\sqrt{\dfrac{\mathrm{mg}}{\pi\rho}}$ for equilibrium}
\opt{D}{Power of engine $=\dfrac{\mathrm{mg}}{r}\sqrt{\dfrac{\mathrm{mg}}{\pi\rho}}$}
\end{qoptions}
\qn{12} Incompressible liquid of negligible viscosity is pumped at total volume flow rate Q through two small holes into the narrow gap between closely spaced parallel plates as shown. The liquid flowing away from the holes has only radial motion. Assume uniform flow across any vertical section.\par
\qfig{p1_phy_s2_q12_1}
\begin{qoptions}
\opt{A}{Flow velocity is proportional to $\dfrac1{r^2}$}
\opt{B}{Flow velocity is proportional to $\dfrac1r$}
\opt{C}{Pressure difference between the centre and any radial point is proportional to $\dfrac1{r^2}$}
\opt{D}{Pressure difference between the centre and any radial point is zero}
\end{qoptions}
\sect{SECTION-II}
\qn{1} A uniform rod of length $\ell$ cross-sectional area A and Young's modulus Y is subjected to two forces `F' and `2F' at its two ends stretching the rod. The heat generated in the rod is found to be $\dfrac{(10\mathrm a+\mathrm b)\,\mathrm F^2\ell}{12\mathrm{AY}}$. Find a+b. \{a, b: Positive integers\}
\qn{2} A rope is hanging vertically with its two ends fixed to walls at same horizontal level (as shown). Tension at the bottom most point is 90 N. If the mass/unit length of rope is 4kg/m, find radius of curvature (in metres) at point `C' \{Take g = 10 m/s$^2$\}\par
\qfig{p1_phy_s3_q02_1}
\qn{3} A source of frequency 1800 Hz is moving on a line, with speed 30 m/s. An observer is standing at a point far distant from the line. The wavelength of sound received by the observer when the source is nearest to the observer is n/11 metres. Find `n' (speed of sound in air =330 m/s)
\qn{4} Ran drops of some initial mass start falling rest. In their way they gain mass from clouds with a rate proportional to the product of instantaneous mass and instantaneous velocity with proportionality constant being 2.5 SI units. Find the terminal speed attained by drops (in m/s). (Take g = 10m/s$^2$)
\qn{5} A sinusoidal wave of wavelength `$\lambda$', amplitude `A' and angular frequency `$\omega$' travel on a string of mass/unit length `$\mu$'. The energy of wave in one wavelength is equal to $\dfrac{\mu\omega^2\mathrm A^2\lambda}{\mathrm b}$. Find `b'
\qn{6} One mole of a monatomic ideal gas follows a process A $\rightarrow$ B, as shown. The specific heat of the process is $\dfrac{13\mathrm R}{6}$. The value of `n' on P-axis (see following figure) is\par
\qfig{p1_phy_s3_q06_1}
\raggedcolumns
\papersec{1}
\partsec{2}{CHEMISTRY}
\sect{SECTION-I (i)}
\qn{1} In which one of the following reactions, products given is correct
\begin{qoptions}
\opt{A}{\qfigi{p1_chem_s1_q01_A}}
\opt{B}{\qfigi{p1_chem_s1_q01_B}}
\opt{C}{\qfigi{p1_chem_s1_q01_C}}
\opt{D}{\qfigi{p1_chem_s1_q01_D}}
\end{qoptions}
\qn{2} The structural isomer of given structure is\par
\qfig{p1_chem_s1_q02_1}
\begin{qoptions}
\opt{A}{\qfigi{p1_chem_s1_q02_A}}
\opt{B}{\qfigi{p1_chem_s1_q02_B}}
\opt{C}{\qfigi{p1_chem_s1_q02_C}}
\opt{D}{\qfigi{p1_chem_s1_q02_D}}
\end{qoptions}
\qn{3} What is the change in Gibbs free energy when 2 moles of an ideal gas at 273K and 2 bar expanded irreversibly to 1 bar
\begin{qoptions}
\opt{A}{$-$ 3147J}
\opt{B}{+3147J}
\opt{C}{$-$2270J}
\opt{D}{+2270J}
\end{qoptions}
\qn{4} To an alkaline salt solution KCN added. Now the solution acidified and FeCl$_3$ added, then red colouration observed. The salt may posses
\begin{qoptions}
\opt{A}{SO$_3^{2-}$}
\opt{B}{S$_2$O$_3^{2-}$}
\opt{C}{NO$_2^-$}
\opt{D}{SO$_4^{2-}$}
\end{qoptions}
\sect{SECTION-I (ii)}
\qn{5} Which of the following is correct ?
\begin{qoptions}
\opt{A}{\qfigi{p1_chem_s2_q05_A}}
\opt{B}{Et$_2$NH $>$ Me$_2$NH (Basic Strength)}
\opt{C}{Et$_3$N $>$ Me$_3$N (Basic Strength)}
\opt{D}{EtNH$_2$ $>$ MeNH$_2$ (Basic Strength)}
\end{qoptions}
\qn{6} Which of the complexes have zero dipole moment, zero magnetic moment and non-zero crystal field stabilization energy (cfse)?
\begin{qoptions}
\opt{A}{trans $-$ [Ni(en)$_2$C$\ell_2$]}
\opt{B}{trans $-$ [Co(NH$_3$)$_4$C$\ell_2$]$^+$}
\opt{C}{trans $-$ [Fe(CN)$_4$(H$_2$O)$_2$]$^{2-}$}
\opt{D}{trans $-$ [Mn(H$_2$O)$_2$(C$\ell$)$_4$]$^{2-}$}
\end{qoptions}
\qn{7}\par
\qfig{p1_chem_s2_q07_1}
In the above reaction which of the following is/are more reactive than benzene diazonium chloride
\begin{qoptions}
\opt{A}{\qfigi{p1_chem_s2_q07_A}}
\opt{B}{\qfigi{p1_chem_s2_q07_B}}
\opt{C}{\qfigi{p1_chem_s2_q07_C}}
\opt{D}{\qfigi{p1_chem_s2_q07_D}}
\end{qoptions}
\qn{8}\par
\qfig{p1_chem_s2_q08_1}
(Given, migratory order Aryl $>$ Arkyl $>$ H)\par
Find, the correct statement(s)
\begin{qoptions}
\opt{A}{[T] can be [Ag(NH$_3$)$_2$ ]OH}
\opt{B}{[T] can be fehling solution}
\opt{C}{[T] can be Cl$_2$ + NaOH}
\opt{D}{[T] can be Benedict solution}
\end{qoptions}
\qn{9} During the hydrolysis of P$_4$O$_{10}$ ,which of the following acids formed as intermediate/ final product
\begin{qoptions}
\opt{A}{Tetrametaphosphoric acid}
\opt{B}{Tetrapolyphosphoric acid}
\opt{C}{Pyrophosphoric acid}
\opt{D}{Phosphoric acid}
\end{qoptions}
\qn{10} When Z $-$ axis is inter nuclear axis which of the following linear combinations are allowed (i.e. give B.M.O \& A.B.M.O)
\begin{qoptions}
\opt{A}{s $-$ p$_\mathrm Z$}
\opt{B}{d$_{\mathrm{XZ}}$ $-$ p$_\mathrm X$}
\opt{C}{d$_{\mathrm{X^2-Y^2}}$ $-$ d$_{\mathrm{X^2-Y^2}}$}
\opt{D}{d$_{\mathrm{YZ}}$ $-$ P$_\mathrm Z$}
\end{qoptions}
\qn{11} Consider an element A crystallizing in simple cubic, face centered cubic and body centered cubic lattice. If the radius of an atom of A is r then which of the following statements is / are correct? (assume no change of radius during formation of different crystals)
\begin{qoptions}
\opt{A}{Shortest distance between two non $-$touching spheres in simple cubic unit cell is $\sqrt8$r}
\opt{B}{Shortest distance between two non $-$ touching spheres in FCC unit cell and SC unit cell is same}
\opt{C}{Shortest distance between two non $-$ touching spheres in BCC unit cell is less than that in FCC unit cell}
\opt{D}{Void space in BCC unit cell is greater than that in FCC unit cell}
\end{qoptions}
\qn{12} Which of the following graphs regarding adsorption are correct?
\begin{qoptions}
\opt{A}{\qfigi{p1_chem_s2_q12_A}}
\opt{B}{\qfigi{p1_chem_s2_q12_B}}
\opt{C}{\qfigi{p1_chem_s2_q12_C}}
\opt{D}{\qfigi{p1_chem_s2_q12_D}}
\end{qoptions}
\sect{SECTION-II}
\qn{1} How many acetals are formed when acetaldehyde is treated with CH$_2$(OH)CH(OH)CH$_2$(OH) in gaseous HCl?
\qn{2} Two liquids A and B form a non $-$ ideal solution which obey the equation P$_\mathrm T$ = P$_\mathrm A^0$ + 3$\left(\mathrm P_B^0-\mathrm P_A^0\right)$ x$_\mathrm B$ + 2$\left(\mathrm P_A^0-\mathrm P_B^0\right)$ x$_\mathrm B^2$. When equi molar mixture of A and B distilled find the mole fraction of B, when this mixture will have a single boiling point.
\qn{3} A 2.0 L flask, initially containing one mole of each CO and H$_2$O, was sealed and heated to 700 K, where the following equilibrium was established CO (g) + H$_2$O$_{(\mathrm g)}$ $\rightleftharpoons$ CO$_2$ (g) + H$_2$ (g); K$_\mathrm C$ = 9. Now the flask was connected to another flask containing some pure CO$_2$(g) at same temperature and pressure, by means of a narrow tube of negligible volume. When the equilibrium was restored, moles of CO was found to be double of its mole at first equilibrium. Determine volume of CO$_2$(g) flask in litres
\qn{4} A sample of H atoms containing all the atoms is some excited state is irradiated with the light of wavelength $\lambda$ and atoms get re $-$ excited to further higher excited state. On removing the light source, atoms de $-$ excite to the ground state with the emission of radiations. The emission spectrum thus obtained consists of ten spectral lines. Out of these ten spectral lines, seven lines have wavelengths smaller and only two line have wavelengths longer than $\lambda$ . The orbit number for initial excited state is
\qn{5} The oxidation number of Cr in the product of alkaline oxidative fusion of FeCr$_2$O$_4$ is
\qn{6} The number of lone pairs of `Xe' present in equatorial position in [XeF$_5$]$^-$
\raggedcolumns
\papersec{1}
\partsec{3}{MATHEMATICS}
\sect{SECTION-I (i)}
\qn{1} If x$\sin$A + y $\sin$B + z$\sin$C = 0 and x$^2\sin$2A + y$^2\sin$2B + z$^2\sin$2C = 0 then the value of x$^3\sin$3A + y$^3\sin$3B + z$^3\sin$3C is equal to (where A,B,C; x, y, z $\in$ R and A + B + C = $\pi$ )
\begin{qoptions}
\opt{A}{0}
\opt{B}{Depend on A, B, C}
\opt{C}{Depend on x, y, z}
\opt{D}{3}
\end{qoptions}
\qn{2} If g (x) is non-negative continuous function such that $\mathrm g(\mathrm x)+\mathrm g\left(\mathrm x+\dfrac12\right)=2$, $\forall$x$\in$R, then the value of $\displaystyle\int_0^1(\sin4\pi\mathrm x)\,\mathrm g(\mathrm x)\,\mathrm{dx}$ is
\begin{qoptions}
\opt{A}{0}
\opt{B}{1}
\opt{C}{$\dfrac12$}
\opt{D}{$\dfrac32$}
\end{qoptions}
\qn{3} If there exists atleast one z satisfying $|z-k\mathrm i|=k+4$ and $|z-3|\le4$, (k $\in$ R) and A = \{0, 4, 8, 10\} then how many values k will take from the set A ?
\begin{qoptions}
\opt{A}{0}
\opt{B}{1}
\opt{C}{2}
\opt{D}{4}
\end{qoptions}
\qn{4} Let $\bar{\mathrm p},\bar{\mathrm q},\bar{\mathrm r}$ are three vectors such that $\bar{\mathrm p}\times\bar{\mathrm r}=\bar{\mathrm q}+\lambda\bar{\mathrm p}$, $\bar{\mathrm p}.\bar{\mathrm r}=1$ then which of the following is true
\begin{qoptions}
\opt{A}{$2\lambda=\dfrac{-\bar{\mathrm p}.\bar{\mathrm q}}{|\bar{\mathrm p}|^2}$}
\opt{B}{$\bar{\mathrm r}=\dfrac{\bar{\mathrm p}-(\bar{\mathrm p}\times\bar{\mathrm q})}{|\bar{\mathrm p}|^2}$}
\opt{C}{$[\bar{\mathrm p}\bar{\mathrm q}\bar{\mathrm r}]=\dfrac{-2|\bar{\mathrm p}\times\bar{\mathrm q}|^2}{|\bar{\mathrm p}|}$}
\opt{D}{$\bar{\mathrm p}\times\bar{\mathrm q}=\bar 0$}
\end{qoptions}
\sect{SECTION-I (ii)}
\qn{5} If P be a point an ellipse $\dfrac{\mathrm x^2}{\mathrm a^2}+\dfrac{\mathrm y^2}{\mathrm b^2}=1$ where foci S, S$^1$ and tangent at P is at a distance d from origin then SP.S$^1$P.d$^2$
\begin{qoptions}
\opt{A}{Independent of `a ' and `b'}
\opt{B}{Independent of ratio $\left(\dfrac{\mathrm a}{\mathrm b}\right)$}
\opt{C}{Independent of eccentric angle of point}
\opt{D}{Less than or equal to max (a$^4$, b$^4$)}
\end{qoptions}
\qn{6} Consider the planes P$_1$ : 2x + y + z + 4 = 0, P$_2$ : y $-$ z + 4 = 0 and P$_3$ : 3x + 2y + z + 8 = 0. Let L$_1$,L$_2$,L$_3$ be the lines of intersection of the planes P$_2$ and P$_3$, P$_3$ and P$_1$, and P$_1$ and P$_2$ respectively. Then which of the following is/are true?
\begin{qoptions}
\opt{A}{Atleast two of the lines L$_1$, L$_2$ and L$_3$ are non-parallel}
\opt{B}{Atleast two of the lines L$_1$, L$_2$ and L$_3$ are parallel}
\opt{C}{The three planes intersect in a line}
\opt{D}{The three planes form a triangular prism.}
\end{qoptions}
\qn{7} f(x) is the function such that $\displaystyle\lim_{x\to0}\dfrac{\mathrm f(\mathrm x)}{\mathrm x}=1$ if $\displaystyle\lim_{x\to0}\dfrac{\mathrm x(1+\mathrm a\cos\mathrm x)-\mathrm b\sin\mathrm x}{(\mathrm f(\mathrm x))^3}=1$ then
\begin{qoptions}
\opt{A}{a = $\dfrac{-5}{2}$}
\opt{B}{b = $\dfrac{-3}{2}$}
\opt{C}{a = $\dfrac{-7}{2}$}
\opt{D}{b = $\dfrac{-5}{2}$}
\end{qoptions}
\qn{8} A die is rolled three times. Let E$_1$ denote the event of getting a number larger than the previous number each time and E$_2$ denote the event that the numbers (in order) form an increasing A.P, then
\begin{qoptions}
\opt{A}{P(E$_2$) $\le$ P(E$_1$)}
\opt{B}{P(E$_1\cap$E$_2$) = $\dfrac1{36}$}
\opt{C}{P (E$_1$) = $\dfrac{10}{3}$P(E$_2$)}
\opt{D}{P (E$_1$) = $\dfrac13$}
\end{qoptions}
\qn{9} If the curve\par\resizebox{\linewidth}{!}{$\mathrm f(\mathrm x)=\displaystyle\int\mathrm x^{27}\left(1+\mathrm x+\mathrm x^2\right)^6\left(6\mathrm x^2+5\mathrm x+4\right)\mathrm{dx}$}\par passes through (0,0), then
\begin{qoptions}
\opt{A}{f (1) = $\dfrac{3^7}{7}$}
\opt{B}{f (1) = $\dfrac{2^7}{7}$}
\opt{C}{f ($-$1) = $\dfrac17$}
\opt{D}{f ($-$1) = $\dfrac{3^7}{14}$}
\end{qoptions}
\qn{10} Normals drawn at points P,Q and R lying on the parabola y$^2$ = 4x intersect at (3,0) then
\begin{qoptions}
\opt{A}{Centroid of $\Delta$PQR = $\left(\dfrac32,0\right)$}
\opt{B}{Circum centre of $\Delta$PQR = $\left(\dfrac53,0\right)$}
\opt{C}{Area of $\Delta$PQR = 2}
\opt{D}{Radius of circum circle of $\Delta$PQR is $\dfrac52$}
\end{qoptions}
\qn{11} Let f(x) be a cubic polynomial which is having local maximum at (1,2) and f$'$(x) has local extrema at x = 0. If f(0) = 1 then
\begin{qoptions}
\opt{A}{Number of real root of f(x) $-$ f($-$x) = 6x $-$ 10 is 1}
\opt{B}{If g(x) + f(x) = 0 then complete interval in which g(x) is increasing is x $\in$ R $-$ ($-$1, 1)}
\opt{C}{The range of k for which f(x) = k has three distinct solutions is k $\in$ (0, 2)}
\opt{D}{Range of the function P($\theta$) = f(sin$\theta$) where $\theta\in$ R is [0, 2]}
\end{qoptions}
\qn{12} If $\sin^2$ x $-$ a$\sin$ x + b = 0 has only one solution in (0, $\pi$) then which of the following statements are correct?
\begin{qoptions}
\opt{A}{a $\in$ ($-\infty$, 1] $\cup$ [2, $\infty$)}
\opt{B}{b $\in$ ($-\infty$, 0] $\cup$ [1, $\infty$)}
\opt{C}{a = 1 + b}
\opt{D}{a = b}
\end{qoptions}
\sect{SECTION-II}
\qn{1} Let $\Delta$ABC right angle at A is inscribed in hyperbola xy = 4. If slope of BC is 2 and coordinates of A are ($\alpha$, $\beta$), then the value of $\dfrac{\alpha^2}{3\beta^2}$ is equal to
\qn{2} A dice is randomly rolled five times the probability that outcome of fifth throw is atleast one of the outcomes of first 4 throws is $\dfrac{\mathrm{pqr}}{6^4}$, then the value of p + q + r is (where pqr is a three digit number)
\qn{3} Let A,B be non-singular matrices of order 3 such that det(A) = 5 and A$^{-1}$B$^2$ + AB = 0, then the value of det (A$^6$ $-$ 2A$^4$B + A$^2$B$^2$) is equal to 10$^\mathrm k$ where k is equal to
\qn{4} The area of the quadrilateral, whose vertices are the root of the equation 1 + z $-$ z$^3$ $-$ z$^4$ = 0, is
\qn{5} Let L$_1$ : 3x + 4y $-$1 = 0 and L$_2$ : 5x $-$ 12 y + 2 = 0 be given two lines let image of every point on L$_1$ w.r.t line L lies on L$_2$ then one of the possible equation of L can be ax + by $-$ 23 = 0, a, b $\in$ I, where $\dfrac{\mathrm b}{\mathrm a}$ is equal to
\qn{6} Let AB be any chord of the circle x$^2$ + y$^2$ $-$ 2x $-$ 6y $-$ 6 = 0 which subtends right angle at the point (2,4) then the locus of the midpoint of AB is x$^2$ + y$^2$ + ax + by + c = 0 where $|$a + b + c$|$ is equal to
\raggedcolumns
\papersec{2}
\partsec{1}{PHYSICS}
\sect{SECTION-I (i)}
\qn{1} Figure shows small part of circuit. [Given C$_1$ = C$_2$ = C$_3$ = C$_4$ = 5 $\mu$F ] which of the following statement is/are TRUE? (Assume steady state)\par
\qfig{p2_phy_s1_q01_1}
\begin{qoptions}
\opt{A}{Potential of point O can be determined but that of B can't be determined}
\opt{B}{Potential of both point O \& B can't be determined}
\opt{C}{If charge on capacitor C$_2$ were also specified, potential of point O can be determined but that of B can't be determined}
\opt{D}{If charge on capacitor C$_2$ were also specified, potential of both the point O \& B can be determined}
\end{qoptions}
\qn{2} A conducting loop in the shape of a right angled isosceles triangle of height 10 cm is kept such that the 90$^\circ$ vertex is very close to an infinitely long conducting wire (see the figure). The wire is electrically insulated from the loop. The hypotenuse of the triangle is parallel to the wire. The current in the triangular loop is in counterclockwise direction and increased at a constant rate of 10 As$^{-1}$. Which of the following statement(s) is (are) true?\par
\qfig{p2_phy_s1_q02_1}
\begin{qoptions}
\opt{A}{There is a repulsive force between the wire and the loop}
\opt{B}{If the loop is rotated at a constant angular speed about the wire, an additional emf of $\left(\dfrac{\mu_o}{\pi}\right)$ volt is induced in the wire.}
\opt{C}{The magnitude of induced emf in the wire is $\left(\dfrac{\mu_o}{\pi}\right)$ volts.}
\opt{D}{The induced current in the wire is in opposite direction to the current along the hypotenuse.}
\end{qoptions}
\qn{3} Assume that the earth changes its shape and turns into an infinite cylinder whose radius and distance of moon from the axis of cylindrical earth remains unchanged and also moon remains spherical. Which of the following statement(s) is/are CORRECT?
\begin{qoptions}
\opt{A}{Gravitational field at the moon due to earth increases}
\opt{B}{Gravitational field at the moon due to earth decreases}
\opt{C}{Speed of the moon in its orbit around the cylindrical earth increases}
\opt{D}{Speed of the moon in its orbit around the cylindrical earth decreases}
\end{qoptions}
\qn{4} An air-filled parallel plate capacitor with the plate area A is connected to a battery with an emf V volt and negligible internal resistance. One of the plates vibrates so that the distance between plates varies as d = d$_0$ + $\alpha\cos\omega$t (where $\alpha\ll$ d$_0$ ). The capacitor breaks down when the instantaneous current reaches to the value `I'. Choose the CORRECT option(s)\par
\qfig{p2_phy_s1_q04_1}
\begin{qoptions}
\opt{A}{The maximum possible amplitude vibrations `a' is $\dfrac{\mathrm{Id}_0^2}{\mathrm{VA}\omega\varepsilon_0}$}
\opt{B}{The maximum possible amplitude vibrations `a' is $\dfrac{2\mathrm{Id}_0^2}{\mathrm{VA}\omega\varepsilon_0}$}
\opt{C}{The instantaneous current in circuit is $\dfrac{\varepsilon_0\mathrm{VA}\omega\alpha}{(\mathrm d_0+\alpha\cos\omega\mathrm t)^2}\sin\omega\mathrm t$}
\opt{D}{The instantaneous current in circuit is $\dfrac{\varepsilon_0\mathrm{VA}\omega\alpha}{2(\mathrm d_0+\alpha\cos\omega\mathrm t)^2}\sin\omega\mathrm t$}
\end{qoptions}
\qn{5} A particle of charge +q and mass `m' moving under the influence of a uniform electric field E$\hat{\mathrm i}$ and uniform magnetic field B$\hat{\mathrm k}$ follows a trajectory from P to Q as shown in figure. The velocities at P and Q are v$\hat{\mathrm i}$ and $-2\hat{\mathrm j}$. Which of the following statement(s) is/are correct?\par
\qfig{p2_phy_s1_q05_1}
\begin{qoptions}
\opt{A}{E = $\dfrac34\left[\dfrac{\mathrm{mv}^2}{\mathrm{qa}}\right]$}
\opt{B}{Rate of work done by the electric field at P is $\dfrac34\left[\dfrac{\mathrm{mv}^3}{\mathrm a}\right]$}
\opt{C}{Rate of work done by the electric field at P is zero}
\opt{D}{Rate of work done by both the fields at Q is zero}
\end{qoptions}
\qn{6} A cavity is taken out from a uniform conducting sphere. Inside the cavity a dipole is placed as shown in the figure. Find which of the following is/are true\par
(Given: q = 10$^{-1}$C, $\ell$ = 0.1 mm, $\theta$ = 30$^\circ$, d = 10 cm, R =12 cm)\par
\qfig{p2_phy_s1_q06_1}
\begin{qoptions}
\opt{A}{Potential at points A \& B are same}
\opt{B}{Potential at points A \& B are different}
\opt{C}{Potential at point P is 0 volt}
\opt{D}{Potential at point P is 2 volt}
\end{qoptions}
\qn{7} Water jet coming out of a stationary horizontal tube at speed v = 2m/s strikes horizontally a massive wall moving in opposite direction with same speed. Water comes to rest relative to wall after striking. Cross-section of jet is 3 cm$^2$ and density of water is 1 gm/cc. Choose the CORRECT statement(s)
\begin{qoptions}
\opt{A}{Force exerted on the wall is 1.2 N}
\opt{B}{Force exerted on the wall is 2.4 N}
\opt{C}{Rate of change of kinetic energy of water jet striking the wall is 19.2 J/s}
\opt{D}{Rate of change of kinetic energy of water jet striking the wall is zero}
\end{qoptions}
\qn{8} A potentiometer as shown in figure in which three batteries between A \& B are connected when the point P is the null point, measured by the potentiometer arrangement. Then choose correct statement(s)\par
\qfig{p2_phy_s1_q08_1}
\begin{qoptions}
\opt{A}{Current passing through 6 V battery is 4A}
\opt{B}{Current passing through 2 V battery is 4A}
\opt{C}{V$_\mathrm A$ $-$ V$_\mathrm B$ = 2V}
\opt{D}{V$_\mathrm A$ $-$ V$_\mathrm B$ = 6V}
\end{qoptions}
\sect{SECTION-I (ii)}
\paraq{9}{10}
Answer Q.9 and Q.10 by appropriately matching the lists based on the information given in the paragraph.\par
{\footnotesize\setlength{\tabcolsep}{1.5pt}\noindent\begin{tabular}{|c|L{2.3cm}|c|L{4.7cm}|}
\hline
 & \multicolumn{1}{c|}{List-I} & & \multicolumn{1}{c|}{List-II} \\ \hline
(I) & Magnitude of angular momentum of system about O increases after the instant shown & (P) & A dipole is hinged on a long smooth fixed horizontal wire. A charge Q (A) is threaded on wire and is free to move on the wire. The system comprises of the dipole and point charge. Gravity is absent \qfigi[0.75]{p2_phy_s2_para9_10_1} \\ \hline
(II) & The net torque on body A is zero about O at the instant shown & (Q) & A system of two charges $-$Q and 2Q (A) of mass m each are connected by a light rod and kept on smooth level ground. An external uniform electric field E$_0$ = mg/Q is applied as shown. Gravity is present \qfigi[0.75]{p2_phy_s2_para9_10_2} \\ \hline
(III) & The net force on negative charge increases in magnitude after the instant shown & (R) & An electron (A) is revolving about a fixed proton in anticlockwise direction. The system comprises of the electron and proton. Gravity is absent \qfigi[0.75]{p2_phy_s2_para9_10_3} \\ \hline
(IV) & Mechanical energy of the system increases after the instant shown (ignore potential energy of external electrical field) & (S) & System consists of a particle (A) of mass m and charge $-$Q. It is projected from level ground in presence of constant electric field E = mg/Q in vertically upward direction. Gravity is present \qfigi[0.75]{p2_phy_s2_para9_10_4} \\ \hline
 & & (T) & System consists of a uniform ring and particle (A) of mass m and charge $-$Q on its circumference as shown. The ring is in pure rolling on a sufficiently rough horizontal surface. Gravity is present and E$_0$ is external electric field \qfigi[0.75]{p2_phy_s2_para9_10_5} \\ \hline
 & & (U) & None of these \\ \hline
\end{tabular}}\par
\qn{9} Choose correct alternative for the above list
\begin{qoptions}
\opt{A}{I-P,Q,S ; II-P,R,S}
\opt{B}{I-P,Q,R ; II-P,R}
\opt{C}{I-P,Q ; II-P,R,S}
\opt{D}{I-P,Q,R ; II-P,R,S}
\end{qoptions}
\qn{10} Choose correct alternative for the above list
\begin{qoptions}
\opt{A}{III-P, Q,T ; IV-P,Q}
\opt{B}{III-P,Q ; IV-Q}
\opt{C}{III-P,Q,R ; IV-P}
\opt{D}{III-P, Q; IV-P}
\end{qoptions}
\paraq{11}{12}
An object O (real) is placed at focus of an equi-biconvex lens as shown in figure. The refractive index of the lens is $\mu$ = 1.5 and the radius of curvature of either surface of lens is R. The lens is surrounded by air. In each statement of List-I, some changes are made to situation given above and information regarding final image formed as a result is given in List-II. The distance between lens and object is unchanged in all statements of List-I. Match the statements in List-I with resulting image in List-II.\par
\qfig{p2_phy_s2_para11_12_1}
{\footnotesize\setlength{\tabcolsep}{1.5pt}\noindent\begin{tabular}{|c|L{3.1cm}|c|L{3.7cm}|}
\hline
 & \multicolumn{1}{c|}{List-I} & & \multicolumn{1}{c|}{List-II} \\ \hline
(I) & If the refractive index of the lens is doubled (that is, made 2$\mu$), then & (P) & Final image is real \\ \hline
(II) & If the radius of curvature is doubled (that is, made 2R), then & (Q) & Final image is virtual \\ \hline
(III) & If a glass slab of refractive index $\mu$ = 1.5 is introduced between the object and lens as shown, then \qfigi[0.75]{p2_phy_s2_para11_12_III} & (R) & Final image is always enlarged. \\ \hline
(IV) & If the left side of lens is filled with a medium of refractive index $\mu$ = 1.5 as shown, then \qfigi[0.75]{p2_phy_s2_para11_12_IV} & (S) & Final image is of same size as the object \\ \hline
 & & (T) & Final image becomes smaller in size in comparison to size of image before the change was made \\ \hline
 & & (U) & Image formation is not possible. \\ \hline
\end{tabular}}\par
\qn{11} Choose correct options for conditions in I and II
\begin{qoptions}
\opt{A}{I-P; II-P,Q,R}
\opt{B}{I-P,T; II-P,Q,T}
\opt{C}{I-P,T; II- Q,T}
\opt{D}{I-P; II-P,Q,S}
\end{qoptions}
\qn{12} Choose correct options for conditions in III and IV
\begin{qoptions}
\opt{A}{III-Q,T; IV-Q,T}
\opt{B}{III-P,T; IV-Q,T}
\opt{C}{III-P; IV- P,Q,R}
\opt{D}{III-Q,T; IV-P,Q,T}
\end{qoptions}
\sect{SECTION-II}
\qn{1} A source of sound with natural frequency f$_0$ moves uniformly along a straight line. A stationary sound detector is located at a distance of 300 m from the straight line. When the source is closest to the detector, it records a frequency of 3 f$_0$ as received from the source. Find the distance (in `m') of the source from the detector at the instant when it record natural frequency f$_0$ of the source. If your answer is $300\sqrt{\dfrac{\mathrm x}{\mathrm y}}$ where x \& y are reduces to nearest integers from 1 to 9 then fill value (x) + (y).
\qn{2} The drag force on a spherical CO$_2$ bubble in water is given by F$_0$ = 6$\pi\eta$rv, where $\eta$ is the viscosity of the surrounding fluid and r is the radius of the sphere (Their relation is called stoke's law). Bubble can be consider to be almost a mass less sphere of radius r moving up in a liquid with terminal speed. Ignore any expansion of bubble during rise. Calculate time taken by sphere to travel distance 20 cm in seconds. (Given; Radius r = 3 mm, density of water $\rho$ = 10$^3$ kg/m$^3$, Density of CO$_2$ = 1kg/m$^3$, g = 10m/s$^2$ , $\eta$ = 0.5 SI unit)
\qn{3} We would; like to increase the length of a 15 cm long copper rod of cross-section 4 mm$^2$ by 1 mm. The energy absorbed by the rod if it is heated is E$_1$. The energy stored by the rod if it is stretched slowly is E$_2$. Then find E$_1$ / E$_2$. If your answer is x fill value of x/100.\par
[Various parameters of Copper are : Density = 9 $\times$ 10$^3$ kg/m$^3$, Thermal coefficient of linear expansion = 8 $\times$ 10$^{-5}$ K$^{-1}$, Young's modulus = 135 $\times$ 10$^9$ Pa , Specific heat = 400 J / kg $-$ K]
\qn{4} A heat engine takes 0.4 moles of an ideal gas around the cycle as shown in the figure. Process 1 $\rightarrow$ 2 is isochoric, process 2 $\rightarrow$ 3 is adiabatic and process 3 $\rightarrow$ 1 is isobaric. The value of $\gamma$ for the gas is 5/3. Find the net heat flow during the cycle in joule. If your answer is N fill value of N/50 ?\par
[Take: Universal gas constant (R) = 50 / 6 J mol$^{-1}$K$^{-1}$ : 1 atm = 10$^6$ N/m$^2$ and 2$^{3/5}$ = 1.5]\par
\qfig{p2_phy_s3_q04_1}
\qn{5} A certain planet of radius R is composed of a uniform material which generates net power P through radioactive decay. The heat conducted per second $\dfrac{\Delta\mathrm Q}{\Delta\mathrm t}=\mathrm P\dfrac{\mathrm r^3}{\mathrm R^3}$, where r is radial distance from centre. This results in a temperature differential between the inside and outside of the planet as heat is transferred from the interior to the surface. The rate of heat transfer inside is governed by conduction. It is found that thermal conductivity k is constant for the planet. Assume that the planet is in a steady state i.e., temperature depends on position, but does not depend on time. (Assuming black body radiation and emissivity is 1). If T$_\mathrm C$ and T$_\mathrm S$ are temperatures at center and surface respectively. Then temperature difference T$_\mathrm C$ $-$ T$_\mathrm S$ = $\dfrac{\sigma\mathrm R^{\mathrm m}\mathrm T_{\mathrm S}^{\mathrm n}}{\ell\mathrm k}$ fill value of (m + n + $\ell$)
\qn{6} When an $\alpha$-particle collides with a stationary lithium nucleus, the following nuclear reaction may Take place $^4$He+$^7$Li$\rightarrow^{10}$B + n if Q-value of this reaction is E = $-$2.8 MeV, find the minimum Value of kinetic energy (in MeV) of incoming $\alpha$-particles required for the reaction.
\raggedcolumns
\papersec{2}
\partsec{2}{CHEMISTRY}
\sect{SECTION-I (i)}
\qn{1} In polar protic solvent, among the following which is/are strongly nucleophilic than acetate ion
\begin{qoptions}
\opt{A}{OH$^-$}
\opt{B}{C$_6$H$_5$O$^-$}
\opt{C}{C$_6$H$_5$S$^-$}
\opt{D}{\qfigi{p2_chem_s1_q01_D}}
\end{qoptions}
\qn{2} Which of the following in not a correct conformer pair?
\begin{qoptions}
\opt{A}{\qfigi{p2_chem_s1_q02_A}}
\opt{B}{\qfigi{p2_chem_s1_q02_B}}
\opt{C}{\qfigi{p2_chem_s1_q02_C}}
\opt{D}{\qfigi{p2_chem_s1_q02_D}}
\end{qoptions}
\qn{3} In parallel path reactions, product from stable intermediate is called K.C.P and it is major product at low temperature. The stable product is called T.C.P and it is major product at high temperature. In which of the following reactions products mentioned are correct ?
\begin{qoptions}
\opt{A}{\qfigi{p2_chem_s1_q03_A}}
\opt{B}{\qfigi{p2_chem_s1_q03_B}}
\opt{C}{\qfigi{p2_chem_s1_q03_C}}
\opt{D}{\qfigi{p2_chem_s1_q03_D}}
\end{qoptions}
\qn{4} With which of the following reagents, ferrous ions give red colouration in the presence of ferric ions
\begin{qoptions}
\opt{A}{KSCN}
\opt{B}{NH$_3$ + dmg}
\opt{C}{O-Phenanthroline}
\opt{D}{2, 2$'$-Bipyridyl}
\end{qoptions}
\qn{5} Which of the following molecular orbitals is/are gerade
\begin{qoptions}
\opt{A}{$\sigma_{\mathrm{p-p}}$}
\opt{B}{$\sigma^*_{\mathrm{p-p}}$}
\opt{C}{$\delta_{\mathrm{d-d}}$}
\opt{D}{$\pi^*_{\mathrm{p-p}}$}
\end{qoptions}
\qn{6} Which of the following are soluble in CS$_2$ ?
\begin{qoptions}
\opt{A}{White phosphorus}
\opt{B}{Red phosphorus}
\opt{C}{Orthorhombic sulphur}
\opt{D}{Monoclinic sulphur}
\end{qoptions}
\qn{7} During the extraction of copper from copper pyrites, in the Bessemer convertor which of the following reactions take place
\begin{qoptions}
\opt{A}{FeO + SiO$_2$ $\rightarrow$ FeSiO$_3$}
\opt{B}{Cu$_2$S + 2Cu$_2$O $\rightarrow$ Cu + SO$_2$}
\opt{C}{Cu$_2$S + O$_2$ $\rightarrow$ Cu$_2$O + SO$_2$}
\opt{D}{C + O$_2$ $\rightarrow$ CO$_2$}
\end{qoptions}
\qn{8} 10 ml mixture containing 0.1 M NaHCO$_3$ and 0.1M Na$_2$CO$_3$ is titrated against an acid with an acid$-$base indicator. Among the following which observations is/are correct
\begin{qoptions}
\opt{A}{When phenolphthalene indicator used, end point observed when 10ml of 1M HCl is added}
\opt{B}{After getting end point with phenolphthalein if we added methyl orange indicator, to get next end point 20 ml of 0.1M HCl is to be added furtherly}
\opt{C}{When methyl orange indicator used from the beginning, end point observed when 6ml of 0.5M HCl added}
\opt{D}{When phenolphthalene indicator used, gram equivalent weight of Na$_2$CO$_3$ is 106 g but when methyl orange indicator used gram equivalent weight of Na$_2$ CO$_3$ is 53g}
\end{qoptions}
\sect{SECTION-I (ii)}
\paraq{9}{10}
Answer Q.9 and Q.10 by appropriately matching the lists based on the information given in the paragraph.\par
List-I includes starting materials and reagents. List$-$II gives the structures of product compounds from the reactions of List$-$I.\par
{\footnotesize\setlength{\tabcolsep}{1.5pt}\noindent\begin{tabular}{|c|L{4.1cm}|c|L{2.5cm}|}
\hline
 & \multicolumn{1}{c|}{List-I} & & \multicolumn{1}{c|}{List-II} \\ \hline
(I) & \qfigi[0.75]{p2_chem_s2_para9_10_1} & (P) & \qfigi[0.75]{p2_chem_s2_para9_10_5} \\ \hline
(II) & \qfigi[0.75]{p2_chem_s2_para9_10_2} & (Q) & \qfigi[0.75]{p2_chem_s2_para9_10_6} \\ \hline
(III) & \qfigi[0.75]{p2_chem_s2_para9_10_3} & (R) & \qfigi[0.75]{p2_chem_s2_para9_10_7} \\ \hline
(IV) & \qfigi[0.75]{p2_chem_s2_para9_10_4} & (S) & \qfigi[0.75]{p2_chem_s2_para9_10_8} \\ \hline
 & & (T) & \qfigi[0.75]{p2_chem_s2_para9_10_9} \\ \hline
 & & (U) & \qfigi[0.75]{p2_chem_s2_para9_10_10} \\ \hline
\end{tabular}}\par
\qn{9} Which of the following has correct combination considering List $-$ I and List $-$ II ?
\begin{qoptions}
\opt{A}{I-P}
\opt{B}{I-Q}
\opt{C}{II-P}
\opt{D}{III-S}
\end{qoptions}
\qn{10} Which of the following has correct combination considering List $-$ I and List $-$ II ?
\begin{qoptions}
\opt{A}{III $-$ R}
\opt{B}{III $-$ S}
\opt{C}{IV $-$ T}
\opt{D}{II $-$ P}
\end{qoptions}
\paraq{11}{12}
Answer Q.11 and Q.12 by appropriately matching the lists based on the information given in the paragraph.\par
List$-$I includes various thermodynamic processes involved for n moles of ideal/real gas in the expansion from P$_1$, V$_1$, T$_1$ to V$_2$. List$-$II includes change in various thermodynamic properties for the processes mentioned in List$-$I\par
{\footnotesize\setlength{\tabcolsep}{1.5pt}\noindent\begin{tabular}{|c|L{3.6cm}|c|L{3.2cm}|}
\hline
 & \multicolumn{1}{c|}{List-I} & & \multicolumn{1}{c|}{List-II} \\ \hline
(I) & Isothermal reversible expansion of ideal gas & (P) & $\Delta\mathrm S_{\mathrm{sys}}=\mathrm{nR}\ln\left(\dfrac{\mathrm V_2}{\mathrm V_1}\right)$ \\ \hline
(II) & Isothermal irreversible expansion of ideal gas & (Q) & $\Delta\mathrm S_{\mathrm{total}}=0$ \\ \hline
(III) & Adiabatic free expansion of ideal gas & (R) & $\Delta\mathrm S_{\mathrm{total}}$ is + ve \\ \hline
(IV) & Adiabatic free expansion of real gas & (S) & $\Delta\mathrm U=0$ \\ \hline
 & & (T) & $\Delta\mathrm T\neq0$ \\ \hline
 & & (U) & $\Delta\mathrm S_{\mathrm{total}}=\Delta\mathrm S_{\mathrm{sys}}$ \\ \hline
\end{tabular}}\par
\qn{11} Among the following incorrect match is
\begin{qoptions}
\opt{A}{II-P}
\opt{B}{I-Q}
\opt{C}{III-U}
\opt{D}{III-T}
\end{qoptions}
\qn{12} Among the following correct match is
\begin{qoptions}
\opt{A}{IV-T}
\opt{B}{IV-P}
\opt{C}{III-Q}
\opt{D}{II-Q}
\end{qoptions}
\sect{SECTION-II}
\qn{1} 1.34g of organic compound with the molecular weight 134 gives 7.05g of AgI precipitate in the Ziesels method of estimation of methoxy groups. What is the weight percentage of methoxy groups in the organic compound?\par
(At wt of Ag is 108, I is 127)
\qn{2} Total number of sterio isomers possible for the octahedral complex [M(AB)C$_2$de] are (Where : (AB) is unsymmetrical bi dentate ligand; c, d and e are mono dentate ligands)
\qn{3} In the reaction between permanganate and thiosulphate ions in faintly alkaline solutions, rate of disappearance of permanganate ion is how many times the rate of disappearance of thiosulphate ion?
\qn{4} When three moles of helium is mixed with one mole of neon at constant temperature and pressure (25$^0$C,1 atm), the entropy of mixing (in JK$^{-1}$, rounded off to two decimal place) is\underline{\hspace{2cm}}\par
(Given: R = 8.314JK$^{-1}$ mol$^{-1}$)
\qn{5} In the mixture 1.0 $\times$ 10$^{-3}$ M AgNO$_3$ + 0.20 M NH$_3$, concentration of [Ag (NH$_3$)]$^+$ in correct scientific notation is P $\times$ 10$^{-\mathrm P}$M. The value of `P' is\par
Given : [Ag(NH$_3$)$_2$]$^+$ $\rightleftharpoons$ [Ag(NH$_3$)]$^+$ + NH$_3$; K$_1$ = 1.4 $\times$ 10$^{-6}$M\par
[Ag(NH$_3$)]$^+$ $\rightleftharpoons$ Ag$^+$ + NH$_3$ ; K$_2$ = 4.3 $\times$ 10$^{-4}$M
\qn{6} Third virial coefficient for a real gas is 1.6 $\times$ 10$^{-3}$ dm$^6$ / mole$^2$ critical volume of that real gas (in cm$^3$ / mole) is
\raggedcolumns
\papersec{2}
\partsec{3}{MATHEMATICS}
\sect{SECTION-I (i)}
\qn{1} Consider a parallelepiped OAC$'$ BA$'$ CB$'$ O$'$ where OA = 2, OB = 3 and OC = 4, if areas of $\Delta$OPQ, $\Delta$OQR and $\Delta$OPR are least (as shown in figure), then\par
\qfig{p2_math_s1_q01_1}
\begin{qoptions}
\opt{A}{Volume of tetrahedron OPQR is 24 cubic units}
\opt{B}{Volume of tetrahedron OPQR is 32 cubic units}
\opt{C}{Shortest distance between BA$'$ \& PR is 3}
\opt{D}{Shortest distance between BA$'$ \& PR is 6}
\end{qoptions}
\qn{2} If the sum to infinite terms of the series $\dfrac1{1.3.5.7.9}+\dfrac1{3.5.7.9.11}+\dfrac1{5.7.9.11.13}+\ldots.$ is `k', then
\begin{qoptions}
\opt{A}{`k' is rational number}
\opt{B}{k = $\dfrac1{840}$}
\opt{C}{k = $\dfrac1{842}$}
\opt{D}{$\dfrac1{\mathrm k}$ is a multiple of 2}
\end{qoptions}
\qn{3} The equation\par
\resizebox{\linewidth}{!}{$\begin{vmatrix}(1+x)^2&(1+x)^2&-(2+x^3)\\2x+1&3x&1-5x\\x+1&2x&2-3x\end{vmatrix}+\begin{vmatrix}(1+x)^2&2x+1&x+1\\(1-x)^2&3x&2x\\1-2x&3x-2&2x-3\end{vmatrix}=0$}\par
has
\begin{qoptions}
\opt{A}{No real solution}
\opt{B}{Exactly 4 real solutions}
\opt{C}{Exactly two real and two non-real solutions}
\opt{D}{Infinite number of solutions, real or non-real}
\end{qoptions}
\qn{4} Let X = \{1,2,3,5,6\} and a, c are natural numbers selected from set X with replacement.\par
Let $\mathrm f(\mathrm x)=\begin{cases}\mathrm x^2+\mathrm x&\mathrm x\geqslant1\\\mathrm x^2-\mathrm{ax}&\mathrm x>1\end{cases}$, $\mathrm g(\mathrm x)=\begin{cases}\mathrm x^2-\mathrm{cx},&\mathrm x\leqslant2\\\mathrm x^2-4\mathrm x,&\mathrm x>2\end{cases}$ and h(x) = f(x).g(x), x $\in$ R\par
N$_1$ = Number of function h(x) such that f(x) and g(x) are both discontinuous\par
N$_2$ = Number of function h(x) such that h(x) is continuous at x = 1 and discontinuous at x = 2 but f(x) and g(x) are both discontinuous\par
N$_3$ = Number of function h(x) such that h(x) is discontinuous at x = 1 and continuous at x = 2, Then
\begin{qoptions}
\opt{A}{N$_1$= 30}
\opt{B}{N$_2$ = 6}
\opt{C}{N$_3$ = 10}
\opt{D}{N$_1$ + N$_2$ + N$_3$ = 45}
\end{qoptions}
\qn{5} Let $\mathrm g(\mathrm x)=\displaystyle\int_0^{\mathrm x}\left(3\mathrm t^2+2\mathrm t+9\right)\mathrm{dt}$ and f(x) be a decreasing function, $\forall$x $\geqslant$ 0 such that $\overrightarrow{\mathrm{AB}}$ = f(x)$\hat{\mathrm i}$ + g(x)$\hat{\mathrm j}$ and $\overrightarrow{\mathrm{AC}}$ = g(x)$\hat{\mathrm i}$ + f(x)$\hat{\mathrm j}$ are the two smallest sides of $\Delta$ABC whose circumcenter lies outside the triangle, $\forall$x $>$ 0 then
\begin{qoptions}
\opt{A}{f(x) $<$ 0}
\opt{B}{g(x) $<$ 0}
\opt{C}{g(x) $>$ 0}
\opt{D}{$\displaystyle\lim_{t\to0}\lim_{x\to\infty}\left(\cot\left(\dfrac\pi4\left(1-t^2\right)\right)\right)^{f(x)g(x)}=\mathrm e$}
\end{qoptions}
\qn{6} A and B throw alternately throw a pair of dice. A wins if he throws a sum 6 before B throws 7 and B wins if he throws a 7 before A throws 6. If A starts the game, then
\begin{qoptions}
\opt{A}{Probability that A wins is $\dfrac{31}{61}$}
\opt{B}{Probability that B wins is $\dfrac{31}{61}$}
\opt{C}{Probability that A wins is $\dfrac{30}{61}$}
\opt{D}{Probability that B wins is $\dfrac{30}{61}$}
\end{qoptions}
\qn{7} Let f(x) = $\cos^2$x. e$^{\tan \mathrm x}$. $\forall$x $\in\left(\dfrac{-\pi}2,\dfrac\pi2\right)$ then
\begin{qoptions}
\opt{A}{$f'$(x) has a point of local minimum at x = $\dfrac\pi4$}
\opt{B}{$f'$(x) has a point of local maximum in $\left(\dfrac{-\pi}4,0\right)$}
\opt{C}{$f'$(x) has exactly two points of local maximum/minimum in $\left(\dfrac{-\pi}2,\dfrac\pi2\right)$}
\opt{D}{$f''$(x) = 0 has no root in $\left(\dfrac{-\pi}2,\dfrac\pi2\right)$}
\end{qoptions}
\qn{8} The equation 10x$^4$ $-$ 3x$^2$ $-$ 1 = 0 has
\begin{qoptions}
\opt{A}{Atleast one root in ($-$1,0)}
\opt{B}{Atleast one root in (0,1)}
\opt{C}{Atleast two root in ($-$1,1)}
\opt{D}{No root in ($-$1,1)}
\end{qoptions}
\sect{SECTION-I (ii)}
\paraq{9}{10}
Answer Q.9 and Q.10 by appropriately matching the lists based on the information given in the paragraph.\par
{\footnotesize\setlength{\tabcolsep}{1.5pt}\noindent\begin{tabular}{|c|L{5.4cm}|c|c|}
\hline
 & \multicolumn{1}{c|}{List-I} & & List-II \\ \hline
(I) & Let f : R $\longrightarrow$ R and\newline f$_\mathrm n$(x) = f (f$_{\mathrm n-1}$(x))$\forall$n $\geqslant$ 2, n $\in$ N, the roots of equation\newline f$_3$(x)f$_2$(x)f(x) $-$ 25f$_2$(x)f(x) + 175f(x) = 375\newline which also satisfy equation f(x) = x will be & (P) & 1 \\ \hline
(II) & Let f : [5, 10] onto [4,17], the integers in the range of y = f (f (f(x))) is/are & (Q) & 5 \\ \hline
(III) & Let f(x) = 8cot$^{-1}$(cot x) + 5sin$^{-1}$(sin x) + 4tan$^{-1}$(tan x) $-$ sin (sin$^{-1}$x), then possible integral values which f(x) can take & (R) & 10 \\ \hline
(IV) & Let `a' denotes the roots of equation\newline cos(cos$^{-1}$x) + sin$^{-1}$ sin$\left(\dfrac{1+\mathrm x^2}2\right)$ = 2sec$^{-1}$(sec x)\newline then possible values of [$|$10a$|$] where[.] denotes the greatest integer function will be & (S) & 15 \\ \hline
 & & (T) & $-$3 \\ \hline
 & & (U) & 0 \\ \hline
\end{tabular}}\par\vspace{-6pt}
\qn{9} Which of following is correct combination?
\begin{qoptions}
\opt{A}{(I) (P) (R)}
\opt{B}{(I) (Q) (R)}
\opt{C}{(II) (Q) (R) (S)}
\opt{D}{(II) (P) (S)}
\end{qoptions}
\qn{10} Which of following is correct combination?
\begin{qoptions}
\opt{A}{(III) (P) (Q) (T) (U)}
\opt{B}{(III) (P) (Q) (R) (U)}
\opt{C}{(IV) (P) (R)}
\opt{D}{(IV) (Q) (S)}
\end{qoptions}
\paraq{11}{12}
Answer Q.11 and Q.12 by appropriately matching the lists based on the information given in the paragraph\par
Let C be the set of curves having the property that the point of intersection of tangent with y-axis is equidistant the point of tangency and origin (0, 0)\par
{\footnotesize\setlength{\tabcolsep}{1.5pt}\noindent\begin{tabular}{|c|L{4.9cm}|c|L{1.9cm}|}
\hline
 & \multicolumn{1}{c|}{List-I} & & \multicolumn{1}{c|}{List-II} \\ \hline
(I) & If C$_1$, C$_2$ $\in$ C\newline C$_1$ : Curve is passing through (1, 0)\newline C$_2$ : Curve is passing through ($-$1, 0)\newline Then number of common tangents for C$_1$ and C$_2$ & (P) & has a factor 2 \\ \hline
(II) & If C$_3$ $\in$ C, C$_3$ : is passing through (2, 4). If $\dfrac{\mathrm x}{\mathrm a}+\dfrac{\mathrm y}{\mathrm b}=1$ is tangent to C$_3$, then If k$_1$a + k$_2$b$^2$ + k$_3$ab$^2$ = 0, then k$_1$ $-$ k$_2$ + k$_3$ (consider k$_1$ $-$ k$_2$ + k$_3$ $>$ 0) & (Q) & has a factor 3 \\ \hline
(III) & If common tangents of C$_1$ and C$_2$ form an equilateral triangle, where C$_1$, C$_2$ $\in$ C and C$_1$ : curve passes through (2, 0), if distance between centres of C$_1$ \& C$_2$ is K then 3K & (R) & less than 10 \\ \hline
(IV) & If C$_1$, C$_2$ $\in$ C\newline C$_1$ : Curve is passing through (1, 0)\newline C$_2$ : Curve is passing through ($-$1, 0)\newline Distance between the centres of the curves C$_1$ \& C$_2$ & (S) & greater than 5 \\ \hline
 & & (T) & has factor 5 \\ \hline
 & & (U) & has factor 6 \\ \hline
\end{tabular}}\par
\qn{11} Which of following is correct combination ?
\begin{qoptions}
\opt{A}{(I) (Q) (S)}
\opt{B}{(I) (P) (R)}
\opt{C}{(II) (P) (S)}
\opt{D}{(II) (Q) (S)}
\end{qoptions}
\qn{12} Which of following is correct combination ?
\begin{qoptions}
\opt{A}{(III) (P) (R)}
\opt{B}{(III) (P) (S)}
\opt{C}{(IV) (P) (S)}
\opt{D}{(IV) (Q) (S)}
\end{qoptions}
\sect{SECTION-II}
\qn{1} Let f(x) be a differentiable function such that f(f(x)) = x, where x$\in$(0, 1], also f(0) = 1, then if the value of $\displaystyle\int_0^1(\mathrm x-\mathrm f(\mathrm x))^{2020}\mathrm{dx}$ is I then $\dfrac1{\mathrm I}$ is
\qn{2} $\displaystyle\lim_{n\to\infty}\mathrm n^2\int_0^{\frac1{\mathrm n}}\mathrm x^{2019\mathrm x+1}\mathrm{dx}$
\qn{3} On circle $|z-2\mathrm i|=2$ ; through origin `O' a chord OP is drawn. On the tangent at O, a point T is taken so that OT = OP. If TP produced meet normal at O to the circle at Q, then limiting value of OQ as P moves towards origin is\underline{\hspace{1cm}}\par
(Note: P and T lie on the same side of imaginary axis)
\qn{4} Let f(x) = $\cos^{-1}\left(\dfrac{2\mathrm x}{1+\mathrm x^2}\right)$ ,g(x) = $\cot^{-1}\left(\dfrac{2\mathrm x}{\mathrm x^2-1}\right)$ where x $\in$ ($-$1,1). If area bounded by the curves y = f(x) + g(x) and y = $\pi$x$^2$ is A then find [A]\par
[Note: [K] denotes greatest integer less than or equal to K.]
\qn{5} If $\bar{\mathrm a},\bar{\mathrm b}$ are vectors perpendicular to each other and $|\bar{\mathrm a}|=2$, $|\bar{\mathrm b}|=3$, $\bar{\mathrm c}\times\bar{\mathrm a}=\bar{\mathrm b}$, then the least value of 2$|\bar{\mathrm c}-\bar{\mathrm a}|$ is
\qn{6} The shortest distance between origin and a point on the space curve\par
x = 2sint, y = 2cost, z = 3t is
\end{multicols}
\end{document}
